\section{EXPERIMENTAL EVALUATION}

\subsection{Evidence Policy and Artifact Status}
The experimental methodology follows a strict code-to-claim rule: every quantitative manuscript result must be generated from a committed executable configuration and raw per-seed trace. A plot script that directly constructs target curves or hard-codes final percentages is not treated as experimental evidence. Accordingly, legacy numerical curves that cannot be traced to raw simulator output are excluded from the security claims of this revision.

The final artifact contains two explicitly separated evaluation paths. First, the released reduced-order CPS integration backend executes the complete PATBU--risk-gate--trust-gated-adaptation--shield--post-shield-attack composition and is used only for mechanism-level validation. Second, the native MATLAB R2024b path acquires MathWorks' \texttt{IEEE39BusSystem}, discovers candidate observation and control edges, validates the selected actuator edges through finite-difference causal probing, instruments the validated native model, and executes the same seven methods, four corruption cells, and 20 matched seeds. The authoritative native run uses four generator-local physical rotor-speed/\texttt{Pref} pairs, corresponding to Gen10@Bus30, Gen1@Bus39, Gen2@Bus31, and Gen3@Bus32. The resulting $4\times4$ finite-difference sensitivity matrix has rank four and condition number 5.13. Native IEEE 39-bus numerical claims are taken only from the orientation-safe run whose committed raw traces, summary, statistics, and generated LaTeX all share one configuration and provenance record.

\subsection{Benchmarks and Environment}
The primary benchmark is the IEEE 39-bus system. It contains 10 generators and 39 buses and is partitioned into interacting control regions. The existing manuscript also uses the IEEE 118-bus system for scale studies; quantitative 118-bus claims should be retained only when their raw execution traces are made available. Additional IEEE 33-bus or 8500-node results must not be presented as existing results unless corresponding executable adapters and raw logs are added to the artifact.

The legacy experiment description uses a workstation with an Intel Xeon Gold 6230 CPU, 128 GB RAM, an NVIDIA RTX 3090 GPU, Ubuntu 20.04, Python 3.9, PyTorch 2.0, and a Simulink--Python co-simulation stack. The authoritative native IEEE 39-bus evidence instead comes from the released GitHub Actions workflow using MATLAB R2024b with Simulink, Simscape, and Simscape Electrical at commit \texttt{cc315ff522c90ce6adc09df3020472dd9ad820ea}. It executes the native plant at a $1/60$~s sample time, records the validated mapping, configuration hash, seed, and raw trace, and reuses one Fast Restart compile across the $7\times4\times20=560$ experiment matrix through numeric tunable selectors. These CI measurements are reported separately from the legacy workstation description.

\subsection{Methods and Ablations}
The security-focused comparison includes:
\begin{itemize}
    \item a standard federated RL baseline without security controls;
    \item a robust-aggregation FRL baseline;
    \item AEFC without provenance-aware belief processing;
    \item AEFC without the risk-aware recovery gate;
    \item AEFC without trust-gated online adaptation;
    \item AEFC without the physical shield; and
    \item the complete AEFC design.
\end{itemize}
The earlier ``NoCW'' and credibility-weighted ablations are excluded unless credibility weighting is reintroduced into the executable algorithm. Heterogeneous adapters and representation alignment can be evaluated separately as utility-oriented ablations.

\subsection{Joint-Corruption Attack Matrix}
The headline experiment is a $2\times2$ decomposition:
\begin{equation}
\{\text{physical-channel corruption off/on}\}
\times
\{\text{knowledge poisoning off/on}\}.
\end{equation}
This produces clean, knowledge-only, physical-only, and joint-attack conditions. The same scenario names, compared methods, and matched seeds are used in the mechanism-level and native IEEE 39-bus paths so that differences can be attributed to the plant/backend rather than to a changed experimental matrix.

Observation attacks include bounded false-data corruption on the selected native observation channels. Control corruption is injected only after the robust shield and before the validated native actuator edge, preserving the threat-model ordering $u_t^\pi\rightarrow u_t^\star\rightarrow u_t^{\mathrm{exec}}$. Knowledge attacks include model replacement, sign-changing poisoned updates, stale recovery evidence, and colluding Byzantine clients. The native mapping is not inferred from names alone: a control edge is admitted only after a finite-difference probe produces a finite, nontrivial response in the selected native observations.

An additional out-of-envelope stress test should violate the modeled control-attack budget. Failure is acceptable in this setting because the purpose is to demonstrate the theorem boundary:
\begin{quote}
inside the modeled attack envelope, the formal guarantee applies; outside it, no formal safety guarantee is claimed.
\end{quote}

\subsection{Raw-Trace Contract}
Each episode records at least:
\begin{table}[htbp]
\centering
\caption{Minimum trace fields for code-to-claim reproducibility.}
\begin{tabular}{ll}
\toprule
Category & Required fields\\
\midrule
Identity & commit SHA, benchmark, seed, configuration\\
Attack & type, start/end, magnitude, Byzantine ratio\\
State & native observation, normalized observation, safety envelope\\
Belief & belief distortion, clean reference, provenance/trust state\\
Federation & robust aggregate and poisoning condition\\
Policy & deployed-policy drift, adaptation decision\\
Authority & risk bound, gate decision, validation decision\\
Control & shield correction, post-shield attack, executed correction\\
Safety & native safety margin, violation indicator\\
Recovery & recovery state and recovery time\\
\bottomrule
\end{tabular}
\end{table}

Logging the shield correction and post-shield corruption separately is essential because it distinguishes learned-policy authority, shield intervention, and downstream control-channel corruption. In the native model, these signals remain separate until the validated control edge, where the plant receives the original reference plus the bounded shield correction plus the post-shield attack.

\subsection{Security and Recovery Metrics}
The primary safety metrics are the episode-level safety violation rate, the minimum native normalized safety margin, and the shield intervention ratio. Recovery utility is measured using recovery success and time-to-recovery rather than recovery accuracy alone.

Knowledge-integrity metrics include
\begin{equation}
D_b(t)=\|b_t^{\mathrm{attack}}-b_t^{\mathrm{clean}}\|_2,
\end{equation}
\begin{equation}
D_\theta(t)=
\frac{\|\theta_t^{\mathrm{dep}}-\theta_t^{\mathrm{clean}}\|}
{\|\theta_t^{\mathrm{clean}}\|+\epsilon},
\end{equation}
and aggregation deviation
\begin{equation}
D_g(t)=\|\widehat g_t-\mu_t^{H}\|_2.
\end{equation}
The false-authorization rate is
\begin{equation}
\mathrm{FAR}=
\frac{\#\{G_t=1\land
\mathfrak R_t^{\mathrm{true}}>\rho_{\max}\}}
{\#\{\mathfrak R_t^{\mathrm{true}}>\rho_{\max}\}}.
\end{equation}
The shield intervention ratio measures how often the final execution layer changes an authorized learned proposal.

\subsection{Statistical Protocol}
The default evaluation uses 20 matched seeds,
\begin{equation}
\mathrm{seed}\in\{0,1,\ldots,19\},
\end{equation}
with the same seed applied across compared methods in each scenario. Skewed recovery-time and distortion metrics are compared with a paired Wilcoxon signed-rank test. Multiple baseline comparisons use Holm correction. Exact all-zero paired differences are assigned $p=1$ before Holm correction, so a no-effect comparison cannot be converted into an artificial significant result by non-finite-value handling. Effect sizes and 95\% confidence intervals are reported together with $p$-values.

For episode-level safety violation proportions, Wilson confidence intervals are reported. If some episodes do not recover before the simulation horizon, the corresponding recovery-success rate is reported explicitly; time-to-event analysis should be added if censoring is substantial rather than interpreting the horizon as an observed recovery time.

\subsection{Figure and Table Design}
The final paper prioritizes the following evidence:
\begin{enumerate}
    \item a threat/trust-boundary figure showing observation, knowledge, gate, adaptation, shield, downstream control corruption, and plant;
    \item a $2\times2$ joint-corruption comparison;
    \item an attack-episode causal timeline showing $D_b$, $q_t$, $G_t$, $D_\theta$, $u_t^\pi$, $u_t^\star$, $u_t^{\mathrm{exec}}$, and the safety margin;
    \item the distribution of minimum physical safety margins across seeds;
    \item post-poisoning belief/policy correction curves; and
    \item a safety--utility--intervention trade-off over $\rho_{\max}$, $q_{\min}$, and $\Delta_\theta$.
\end{enumerate}

The main results table jointly reports safety violation, minimum safety margin, false authorization, belief distortion, policy drift, recovery success, and recovery time. Communication cost is reported only when computed from actual transmitted bytes in the corresponding trace.

\IfFileExists{results/generated/Results-Auto.tex}{%
\subsection{Mechanism-Level 20-Seed Joint-Corruption Results}

The executable security pipeline was first validated on the released reduced-order CPS backend. This backend exercises the complete PATBU--gate--adaptation--shield path and the post-shield attack ordering, but it is not reported as an IEEE 39-bus result. Each entry below is computed from 20 matched seeds and immutable raw JSONL traces produced by the same commit and configuration.

For the joint physical/knowledge attack, AEFC-Full obtains episode safety-violation rate 0.0000 $\pm$ 0.0000 (95\% normal-approximation CI half-width), compared with 1.0000 $\pm$ 0.0000 for FedRL.

For the joint physical/knowledge attack, AEFC-Full obtains false-authorization rate 0.0000 $\pm$ 0.0000 (95\% normal-approximation CI half-width), compared with 1.0000 $\pm$ 0.0000 for FedRL.

For the joint physical/knowledge attack, AEFC-Full obtains maximum belief distortion 0.4813 $\pm$ 0.0025 (95\% normal-approximation CI half-width), compared with 1.2896 $\pm$ 0.0867 for FedRL.

For the joint physical/knowledge attack, AEFC-Full obtains maximum deployed-policy drift 0.0360 $\pm$ 0.0024 (95\% normal-approximation CI half-width), compared with 3.4224 $\pm$ 0.3630 for FedRL.

For the joint physical/knowledge attack, AEFC-Full obtains recovery success rate 1.0000 $\pm$ 0.0000 (95\% normal-approximation CI half-width), compared with 0.9500 $\pm$ 0.0980 for FedRL.

For the joint physical/knowledge attack, AEFC-Full obtains post-attack recovery time 4.3000 $\pm$ 0.4519 (95\% normal-approximation CI half-width), compared with 40.9000 $\pm$ 2.2161 for FedRL.

These measurements validate executable mechanism composition only. IEEE 39-bus claims remain pending until the licensed R2024b workflow produces native-model traces.
%
}{%
\subsection{Mechanism-Level 20-Seed Results Pending CI}
The manuscript is wired to import the automatically generated mechanism-level results from \texttt{results/generated/Results-Auto.tex}. Until the CI runner has produced the committed summary and immutable raw-trace artifact, no numerical value is inserted here.
}

\IfFileExists{results/ieee39_generated/IEEE39-Results-Auto.tex}{%
\subsection{Native IEEE 39-Bus 20-Seed Joint-Corruption Results}

The native experiment uses the MathWorks R2024b \texttt{IEEE39BusSystem} plant. Four distinct generator-local \texttt{Pref} actuation edges are paired with physical rotor-speed observations from the same generators and accepted only after a full-rank finite-difference causal probe. The same seven methods, four corruption cells, and 20 matched seeds are then executed with control corruption injected downstream of the shield and upstream of the validated native actuator edges.

\begin{table*}[t]
\centering
\caption{Native MathWorks IEEE 39-bus joint-corruption results over 20 matched seeds. Values are means with 95\% confidence intervals; Wilson intervals are used for binary rates.}
\label{tab:ieee39-joint-native}
\scriptsize
\begin{tabular}{lcccccc}
\toprule
Method & SVR $\downarrow$ & FAR $\downarrow$ & Belief dist. $\downarrow$ & Policy drift $\downarrow$ & Rec. success $\uparrow$ & Rec. time $\downarrow$\\
\midrule
FedRL & 0.000 [0.000, 0.161] & 0.000 [0.000, 0.000] & 1.000 [1.000, 1.000] & 1.000 [1.000, 1.000] & 1.000 [0.839, 1.000] & 0.017 [0.017, 0.017]\\
RobustAgg & 0.000 [0.000, 0.161] & 0.000 [0.000, 0.000] & 1.000 [1.000, 1.000] & 1.000 [1.000, 1.000] & 1.000 [0.839, 1.000] & 0.017 [0.017, 0.017]\\
AEFC-no-PATBU & 0.000 [0.000, 0.161] & 0.000 [0.000, 0.000] & 1.000 [1.000, 1.000] & 1.000 [1.000, 1.000] & 1.000 [0.839, 1.000] & 0.017 [0.017, 0.017]\\
AEFC-no-Gate & 0.000 [0.000, 0.161] & 0.258 [0.258, 0.258] & 1.000 [1.000, 1.000] & 1.000 [1.000, 1.000] & 1.000 [0.839, 1.000] & 0.017 [0.017, 0.017]\\
AEFC-no-TGOPA & 0.000 [0.000, 0.161] & 0.258 [0.258, 0.258] & 1.000 [1.000, 1.000] & 1.000 [1.000, 1.000] & 1.000 [0.839, 1.000] & 0.017 [0.017, 0.017]\\
AEFC-no-Shield & 0.000 [0.000, 0.161] & 0.258 [0.258, 0.258] & 1.000 [1.000, 1.000] & 1.000 [1.000, 1.000] & 1.000 [0.839, 1.000] & 0.017 [0.017, 0.017]\\
\textbf{AEFC-Full} & 0.000 [0.000, 0.161] & 0.258 [0.258, 0.258] & 1.000 [1.000, 1.000] & 1.000 [1.000, 1.000] & 1.000 [0.839, 1.000] & 0.017 [0.017, 0.017]\\
\bottomrule
\end{tabular}
\end{table*}


Under joint corruption, AEFC-Full yields episode safety-violation rate 0.0000 (95\% CI [0.0000, 0.1611]), versus 0.0000 ([0.0000, 0.1611]) for FedRL.

Under joint corruption, AEFC-Full yields false-authorization rate 0.2581 (95\% CI [0.2581, 0.2581]), versus 0.0000 ([0.0000, 0.0000]) for FedRL.

Under joint corruption, AEFC-Full yields maximum belief distortion 1.0000 (95\% CI [1.0000, 1.0000]), versus 1.0000 ([1.0000, 1.0000]) for FedRL.

Under joint corruption, AEFC-Full yields maximum deployed-policy drift 1.0000 (95\% CI [1.0000, 1.0000]), versus 1.0000 ([1.0000, 1.0000]) for FedRL.

Under joint corruption, AEFC-Full yields recovery success rate 1.0000 (95\% CI [0.8389, 1.0000]), versus 1.0000 ([0.8389, 1.0000]) for FedRL.

Under joint corruption, AEFC-Full yields post-attack recovery time 0.0167 (95\% CI [0.0167, 0.0167]), versus 0.0167 ([0.0167, 0.0167]) for FedRL.

Continuous integrity, safety-margin, and recovery metrics use paired Wilcoxon tests across matched seeds with Holm correction across the six alternative methods; machine-readable values are released in \texttt{results/ieee39\_generated/ieee39\_joint\_statistics.csv}.
%
}{%
\subsection{Native IEEE 39-Bus Results Pending Licensed CI}
The native R2024b runner and code-to-claim pipeline are implemented, but no IEEE 39-bus numerical result is inserted until the licensed workflow has produced \texttt{interface\_mapping.json}, native per-seed traces, \texttt{results/ieee39\_joint20/summary.csv}, and the automatically generated \texttt{IEEE39-Results-Auto.tex} file.
}
